\section{Dérivation des canaux d'incidence \texorpdfstring{\(90\)}{90} et \texorpdfstring{\(120\)}{120} à partir des demi-couronnes tritiques du Topological Phase Kernel}
\label{sec:sucesor-102-incidencia-tpk-90-120}

Les cardinaux \(90\) et \(120\) sont générés dans la dynamique interne du
Topological Phase Kernel avant toute évaluation exponentielle, réalisation
hexade--octade ou lecture métrologique. Leur provenance est la partition orientée
des \(9\) phases et l'incidence de leurs paires actives ; les canaux
spectraux \(q_\pm\) agiront ultérieurement sur ces objets déjà
construits.

\subsection{Partition tritique de la phase nonadique en demi-couronnes orientées}

Dans la carte canonique \(D_9=\{1,\ldots,9\}\), le sélecteur tritique de phase
détermine la décomposition disjointe
\[
 H_+=\{1,4,7\},\qquad
 H_-=\{2,5,8\},\qquad
 H_0=\{3,6,9\}.
\]
Les classes de \(H_+\) activent le feuillet additif, celles de \(H_-\) le feuillet
multiplicatif et celles de \(H_0\) conservent le seuil. Nous définissons
\[
 H_{\mathrm{act}}=H_+\sqcup H_-,
 \qquad |H_{\mathrm{act}}|=6,
 \qquad \bigl|P_2(H_{\mathrm{act}})\bigr|=\binom62=15,
\]
où \(P_2(H_{\mathrm{act}})\) désigne l'ensemble des paires non ordonnées de
phases actives. Cette partition conserve feuillet, orientation et origine de phase ; ce
n'est pas une classification rétrospective induite par les réalisations physiques.

\subsection{Canal d'incidence de cardinal \texorpdfstring{\(90\)}{90}}

L'incidence interne entre une phase active distinguée et une paire active est
\[
 \mathcal F_{90}
 :=H_{\mathrm{act}}\times P_2(H_{\mathrm{act}}).
\]
Par construction,
\[
 |\mathcal F_{90}|=6\cdot15=90.
\]
Le cardinal \(90\) procède donc de la demi-couronne tritique du TPK. Ce
n'est pas une entrée de Barbero--Immirzi, de la loi des masses ni d'une table
métrologique.

\subsection{Canal d'incidence de cardinal \texorpdfstring{\(120\)}{120}}

En étendant le premier facteur aux \(8\) positions qui précèdent le retour
nonadique, on obtient
\[
 O_{\mathrm{ph}}=\{1,\ldots,8\},
 \qquad
 \mathcal F_{120}
 :=O_{\mathrm{ph}}\times P_2(H_{\mathrm{act}}),
\]
et, par conséquent,
\[
 |\mathcal F_{120}|=8\cdot15=120.
\]
La différence entre \(\mathcal F_{90}\) et \(\mathcal F_{120}\) conserve
l'information de frontière : le second canal intègre les \(8\) positions
antérieures au retour, tandis que le premier conserve uniquement les
\(6\) phases opératoires.

\subsection{Compatibilité dodécaphasique et antériorité causale par rapport à l'évaluation spectrale}

La roue dodécaphasique apparie les \(2\) demi-couronnes orientées et enregistre la
décomposition \(1+5+6=12\), inversible sur son image. Ainsi, les cardinaux
d'incidence sont fixés dans l'ordre causal
\[
 \mathrm{APP}\longrightarrow\mathrm{TRIT}\longrightarrow\mathrm{TPK}
 \longrightarrow(H_+,H_-,H_0)
 \longrightarrow(\mathcal F_{90},\mathcal F_{120}).
\]
Ce n'est qu'ensuite, une fois construites les coordonnées angulaires conjuguées, que les
canaux \(q_\pm\) évaluent spectralement ces incidences au moyen de
\[
 V_{90,120}
 =(I-T^{90})(I-T^{120})^{-1}
 =r_+P_++r_-P_-,
 \qquad
 r_\pm=\frac{1-q_\pm^{90}}{1-q_\pm^{120}}.
\]
L'évaluation \(V_{90,120}\) ne génère pas les cardinaux qui apparaissent dans ses
exposants : elle publie opératoirement les canaux d'incidence produits
auparavant par le TPK.
